\section{Validate}

\paragraph{Input and output.}
Validation consumes a closed AST and returns the same AST when the requested
invariant set holds.  On failure, it reports diagnostics rather than producing
a repaired program.

\paragraph{Contract.}
The pre-flatten validator checks closedness, unresolved-name rejection,
type-alias arity and visibility, and structural properties required before
normalization and flattening.  Later validators check stronger stage-specific
properties such as normalized calls, flat control, backend-required graph shape,
and typed metadata consistency.

\paragraph{Rationale.}
Validation is a certification layer, not a transformation layer.  This is what
lets later stages reject invalid inputs instead of weakening their assumptions.

\paragraph{Formal view.}
A validator is a predicate over an AST representation:
\[
  \mathsf{validate}_k(A) =
  \begin{cases}
  A & \text{if invariant set } k \text{ holds},\\
  \bot & \text{otherwise.}
  \end{cases}
\]
The returned AST is structurally identical to the input.  Diagnostics are
separate from transformation.

\paragraph{Example.}
Before flattening, validation rejects a closed program that still contains a
reachable unresolved name:
\begin{verbatim}
f x = does_not_exist x
\end{verbatim}
After flattening, a stronger validator rejects statement-level control:
\begin{verbatim}
if cond then do
  y <- f x
else do
  y <- g x
\end{verbatim}
because flat Axon must express this through values, helpers, or ternaries.
